% !TeX program = xelatex
% !TeX encoding = UTF-8

\documentclass[11pt,a4paper]{article}

% 中文支持
\usepackage[UTF8]{ctex}

% 页面设置
\usepackage[left=2cm,right=2cm,top=2cm,bottom=2cm]{geometry}
\usepackage{ragged2e}
\usepackage[none]{hyphenat}

% 图形支持
\usepackage{graphicx}
\usepackage{tikz}

% 字体与颜色
\usepackage{xcolor}
\usepackage{fontawesome5}

% 定义主题颜色
\definecolor{primarycolor}{RGB}{41, 128, 185}
\definecolor{secondarycolor}{RGB}{52, 73, 94}
\definecolor{accentcolor}{RGB}{231, 76, 60}
\definecolor{decorcolor}{RGB}{230, 126, 34}

% 设置段落间距
\setlength{\parskip}{0.15em}

\usepackage{titlesec}
\titlespacing{\section}{0pt}{8pt}{5pt}

% 列表格式
\usepackage{enumitem}
\setlist[itemize]{leftmargin=1.2em,topsep=0.1em,itemsep=0.05em}

\begin{document}

% 头部信息区域
\vspace{-0.2cm}
\begin{minipage}[t]{0.75\textwidth}
    \vspace{0.3cm}
    {\Huge \bfseries \color{secondarycolor} 张三}

    \vspace{0.15cm}

    {\large 软件开发工程师}

    \vspace{0.25cm}

    % 联系方式
    \color{secondarycolor}
    {\faPhone} 138-xxxx-xxxx \hspace{0.6cm}
    {\faEnvelope} zhangsan@email.com \hspace{0.6cm}
    {\faMapMarker} 北京市朝阳区

    \vspace{0.1cm}

    {\faWeixin} zhangsan\_wx \hspace{0.6cm}
    {\faGithub} github.com/zhangsan
\end{minipage}%
\hfill
\begin{minipage}[t]{0.2\textwidth}
    \vspace{0.3cm}
    \begin{tikzpicture}
        % 一寸照片框装饰
        \fill[decorcolor!20] (0,0) rectangle (2.5,3.5);
        \draw[primarycolor,thick] (0,0) rectangle (2.5,3.5);
        % 四角装饰
        \foreach \x/\y in {0/0,0/3.5,2.5/0,2.5/3.5}
            \fill[primarycolor] (\x,\y) circle (2pt);
        \node[text width=2cm,align=center,font=\small\color{secondarycolor}] at (1.25,1.75) {一寸照片};
        % 实际使用时，取消下面注释并替换照片路径
        % \node[anchor=south west,inner sep=0] at (0,0) {\includegraphics[width=2.5cm,height=3.5cm]{photo.jpg}};
    \end{tikzpicture}
\end{minipage}

\vspace{0.3cm}

% 教育背景
\section*{\faGraduationCap\quad 教育背景}

\begin{itemize}
    \item \textbf{计算机科学与技术（硕士）} \hfill 2016.09 -- 2019.01

    \textcolor{primarycolor}{清华大学}

    \vspace{0.05em}

    主修课程：高级算法、分布式系统、机器学习、数据库系统、软件工程

    \vspace{0.15em}

    \item \textbf{软件工程（本科）} \hfill 2012.09 -- 2016.06

    \textcolor{primarycolor}{浙江大学}

    \vspace{0.05em}

    GPA: 3.8/4.0，连续三年获得校级奖学金
\end{itemize}

\section*{\faBriefcase\quad 实习/项目经历}

\begin{itemize}
    \item \textbf{高级软件开发工程师} \hfill 2021.06 -- 至今

    \textcolor{primarycolor}{北京科技有限公司}

    \begin{itemize}
        \item 负责公司核心SaaS平台的架构设计与功能开发，用户量突破100万
        \item 主导技术栈从单体架构向微服务架构迁移，系统响应速度提升60\%
        \item 带领5人开发团队，完成多个重要版本迭代，按时交付率达95\%
        \item 引入CI/CD流程，代码部署时间从2小时缩短至15分钟
    \end{itemize}

    \vspace{0.15em}

    \item \textbf{软件开发工程师} \hfill 2019.03 -- 2021.05

    \textcolor{primarycolor}{上海互联网公司}

    \begin{itemize}
        \item 参与电商平台的订单系统和支付系统开发，日均处理订单量50万+
        \item 优化数据库查询性能，将核心接口响应时间从500ms降至100ms
        \item 设计并实现了消息队列系统，有效解决了高并发场景下的消息丢失问题
    \end{itemize}

    \vspace{0.15em}

    \item \textbf{企业级数据分析平台} \hfill 2023.01 -- 2023.08

    \textcolor{primarycolor}{个人项目}

    \begin{itemize}
        \item 使用React + Spring Boot + Kafka构建实时数据分析系统
        \item 设计并实现可视化报表功能，支持自定义报表导出
        \item 系统支持10+TB数据处理，查询响应时间<3秒
    \end{itemize}

    \vspace{0.15em}

    \item \textbf{智能客服系统} \hfill 2022.03 -- 2022.10

    \textcolor{primarycolor}{个人项目}

    \begin{itemize}
        \item 基于NLP技术构建智能问答机器人，解决80\%的用户常见问题
        \item 设计会话管理模块，支持多轮对话和上下文理解
        \item 集成企业微信和钉钉平台，累计服务用户超过50万
    \end{itemize}
\end{itemize}

\section*{\faTools\quad IT技能}

\begin{itemize}
    \item \textbf{编程语言}: Java, Python, JavaScript/TypeScript, Go, SQL
    \item \textbf{前端技术}: React, Vue.js, HTML5/CSS3, Node.js
    \item \textbf{后端框架}: Spring Boot, Django, Flask, gRPC
    \item \textbf{数据库}: MySQL, PostgreSQL, MongoDB, Redis
    \item \textbf{云服务}: AWS, 阿里云, Docker, Kubernetes
    \item \textbf{开发工具}: Git, Jenkins, VS Code, IntelliJ IDEA
\end{itemize}

\section*{\faAward\quad 获奖情况}

\begin{itemize}
    \item \textbf{2022年度优秀员工奖} \hfill 北京科技有限公司
    \item \textbf{PMP项目管理认证} \hfill PMI认证
    \item \textbf{AWS认证架构师} \hfill Amazon Web Services
    \item \textbf{国家级大学生创新创业大赛银奖} \hfill 2015年
\end{itemize}

\section*{\faInfoCircle\quad 其他}

\begin{itemize}
    \item 英语水平：CET-6（520分），能够流畅阅读英文技术文档
    \item 兴趣爱好：摄影、马拉松、开源社区贡献
    \item 自我评价：热爱技术，持续学习，善于沟通与协作
\end{itemize}

\end{document}
